\chapter{校正済み例に用いた式}
\section{線形ばね系}
ばねの微小変位に対し復元力が $-k_{\rm eff}x$ で、質量が$m$の一自由度系では
\[m\ddot{x}+k_{\rm eff}x=0,\qquad \omega_0=\sqrt{k_{\rm eff}/m},\qquad f_0=\omega_0/(2\pi).\]
直列接続では共通の力に対する伸びを加えるため
$1/k_{\rm eff}=\sum_i1/k_i$、並列接続では共通の伸びに対する力を加えるため$k_{\rm eff}=\sum_i k_i$となる。
重力などの定常力が平衡点を移動させても、その点の周囲の線形復元剛性が同じなら固有振動数は変わらない。
